%% Question bank chapter file - generated by bank_tools.py build (do not edit while the bank is still being built:
%% it is regenerated from the per-test data in Question Bank\data). Every question is one bankq block:
%%   \begin{bankq}[id=..., source={...}, paper=.., qno=.., type=..., subtopic={...}, answer={...}, answer_src=key|solved,
%%                 working={...}, difficulty=1-10, diff_parts={K S I C F}, diff_note={...},
%%                 group=<passage id or empty>, also_in={duplicate copies in other tests}]
%% difficulty = min(10, K+S+I+C+F-1): K concepts 1-3, S steps 1-4, I insight 0-2, C computation 0-2, F format 0-1
%%   ...question LaTeX, options, figures (\bfig / \bfigi -> img/)...
%%   \end{bankq}
%% Passages shared by several questions are bankpassage blocks; their questions carry group=<passage id>.
\documentclass[10pt]{article}
\usepackage[a4paper,margin=1.8cm]{geometry}
\usepackage[T1]{fontenc}
\usepackage{amsmath,amssymb,mathtools}
\usepackage{graphicx}
\usepackage[export]{adjustbox}
\usepackage{xcolor,needspace,keyval}
\usepackage{enumitem}
\usepackage{booktabs,array,tabularx}
\usepackage[hidelinks]{hyperref}
\setlength{\parindent}{0pt}
\setlength{\parskip}{3pt}
\makeatletter
\define@key{bq}{id}{\def\bq@id{#1}}\define@key{bq}{source}{\def\bq@src{#1}}\define@key{bq}{paper}{\def\bq@paper{#1}}
\define@key{bq}{qno}{\def\bq@qno{#1}}\define@key{bq}{type}{\def\bq@type{#1}}\define@key{bq}{subtopic}{\def\bq@sub{#1}}
\define@key{bq}{answer}{\def\bq@ans{#1}}\define@key{bq}{answer_src}{\def\bq@asrc{#1}}\define@key{bq}{working}{\def\bq@work{#1}}
\define@key{bq}{group}{\def\bq@group{#1}}\define@key{bq}{also_in}{\def\bq@also{#1}}
\define@key{bq}{difficulty}{\def\bq@diff{#1}}\define@key{bq}{diff_parts}{\def\bq@dparts{#1}}\define@key{bq}{diff_note}{\def\bq@dnote{#1}}
\newenvironment{bankq}[1][]{\def\bq@work{}\def\bq@also{}\def\bq@group{}\def\bq@diff{?}\def\bq@dparts{}\def\bq@dnote{}\setkeys{bq}{#1}%
  \par\needspace{4\baselineskip}\medskip\noindent{\small\textbf{\bq@id}\hfill\color{gray}[\bq@src, Paper \bq@paper, Q\bq@qno]
  \bq@type\ $\cdot$ \bq@sub\ $\cdot$ \textbf{D\bq@diff}}\par\nopagebreak\relax}
  {\par{\small\color{teal}Answer: \textbf{\bq@ans}\ifx\bq@work\@empty\else\ (solved: \bq@work)\fi
  \ifx\bq@also\@empty\else\par\color{gray}Also in: \bq@also\fi}\par}
\define@key{bp}{id}{\def\bp@id{#1}}\define@key{bp}{source}{\def\bp@src{#1}}\define@key{bp}{questions}{\def\bp@qs{#1}}
\newenvironment{bankpassage}[1][]{\setkeys{bp}{#1}\par\needspace{4\baselineskip}\bigskip\noindent
  {\small\textbf{Passage \bp@id}\hfill\color{gray}[\bp@src; questions \bp@qs]}\par\nopagebreak\relax\itshape}{\par}
\makeatother
\newcommand{\bfig}[2]{\par\smallskip{\centering\includegraphics[width=#1pt,max width=0.62\linewidth,max height=0.32\textheight]{img/#2}\par}\smallskip}
\newcommand{\bfigi}[2]{\includegraphics[width=#1pt,max width=\linewidth,max height=0.32\textheight,valign=c]{img/#2}}
\newcommand{\bmissing}[1]{\textit{[figure missing in the source conversion: \detokenize{#1}]}}
\newenvironment{qoptions}{\begin{itemize}[leftmargin=1.8em,itemsep=1pt,topsep=2pt,label={}]}{\end{itemize}}
\newcommand{\opt}[2]{\item[(#1)] #2}
\providecommand{\fwstart}{}\providecommand{\fwend}{}\providecommand{\qfullsplit}{}
\newcolumntype{L}[1]{>{\raggedright\arraybackslash}p{#1}}
\newcolumntype{C}[1]{>{\centering\arraybackslash}p{#1}}
\newcolumntype{R}[1]{>{\raggedleft\arraybackslash}p{#1}}
\providecommand{\sectionhdr}[1]{}\providecommand{\subsectionhdr}[1]{}


\begin{document}
{\Large\bfseries Wave Optics --- Question Bank}\par\smallskip
{\small Physics chapter P20. Questions from all converted test papers, grouped by source series. Answers: printed key where it matches the paper, otherwise solved (answer\_src=solved, with working). Exact repeats across papers appear once, with also\_in listing the other copies.}\par\medskip\hrule

\section*{Score-I papers}
\begin{bankq}[id=SI12-P1-PHY-S1-Q08, source={Score-I - Test 12 (09 Mar 2022)}, paper=1, qno=8, type=multi, subtopic={Interference}, answer={C,D}, answer_src=key, difficulty=4, diff_parts={K1 S2 I1 C0 F1}, diff_note={concepts: interference of two virtual sources | steps: locate the two images; path difference at P; maxima condition for the options | I: images behave as two coherent sources}]
$M_1$ and $M_2$ are two plane mirror \& $S$ is monochromatic source. $AB$ is a `stop' which stops direct light to reach $CF$ and allows reflected light from $M_1$ and $M_2$ to reach $CF$. $P$ is a point on $CF$ such that $SP$ is parallel to $M_1$ \& $M_2$ and perpendicular to $CF$:
\bfig{131.7}{SI12_p1_phy_s1_q08_1.png}
\begin{qoptions}
\opt{A}{Circular fringes will be formed on $CF$}
\opt{B}{$P$ will be central maxima}
\opt{C}{$P$ will be a point of maxima if $\lambda=\dfrac{6d^2}{D}$}
\opt{D}{$P$ will be a point of maxima if $\lambda=\dfrac{d^2}{D}$}
\end{qoptions}
\end{bankq}

\section*{TOAS papers}
\begin{bankq}[id=TOAS17-P1-PHY-S1-Q03, source={TOAS - Test 17 (18 Sep 2022)}, paper=1, qno=3, type=multi, subtopic={YDSE}, answer={B,C,D}, answer_src=key, difficulty=4, diff_parts={K1 S2 I0 C1 F1}, diff_note={concepts: path differences in a double-slit set-up and intensity | steps: path difference to each point for the given d; phase differences; resultant intensity at each point for the four cases}]
In given setup coherent light of wave length $\lambda$ is used assuming $D\gg d$, select correct alternative [each slit produce same intensity on screen]
\bfig{161.1}{TOAS17_p1_phy_s1_q03_1.png}
\begin{qoptions}
\opt{A}{For $d=\sqrt{2\lambda D}$, intensity at C and A will be maximum but minimum at B.}
\opt{B}{For $d=2\sqrt{\lambda D}$, intensity at A, B and C will be maximum intensity.}
\opt{C}{For $d=\sqrt{\dfrac{\lambda D}2}$, intensity at A, B and C will be $\left(\dfrac59\right)^{\text{th}}$ maximum intensity.}
\opt{D}{For $d=\dfrac{\sqrt{\lambda D}}2$, intensity at A and C will be $\left(\dfrac19\right)^{\text{th}}$ maximum intensity.}
\end{qoptions}
\end{bankq}
\begin{bankq}[id=TOAS17-P1-PHY-S1-Q07, source={TOAS - Test 17 (18 Sep 2022)}, paper=1, qno=7, type=multi, subtopic={Diffraction}, answer={A,B,C,D}, answer_src=key, difficulty=3, diff_parts={K1 S2 I0 C0 F1}, diff_note={concepts: double-slit diffraction pattern | steps: missing orders from d/a; envelope width vs wavelength; number of fringes in the envelope}]
Consider a diffraction pattern of 2 slits as shown in figure. The pattern is formed by 2 slits of width $a=3\mu$m each separated by $d=18\ \mu$m.
\bfig{369.7}{TOAS17_p1_phy_s1_q07_1.png}

\begin{qoptions}
\opt{A}{Here 6$^{\text{th}}$ order maxima of interference can't be seen.}
\opt{B}{If the wavelength is decreased, the central envelope becomes narrower.}
\opt{C}{If the wavelength is decreased, the width of each interference maxima becomes lesser.}
\opt{D}{If the wavelength is decreased, same number of interference maxima will be seen in the central envelope of diffraction.}
\end{qoptions}
\end{bankq}
\begin{bankq}[id=TOAS17-P2-PHY-S1-Q01, source={TOAS - Test 17 (18 Sep 2022)}, paper=2, qno=1, type=single, subtopic={Polarization}, answer={B}, answer_src=key, difficulty=1, diff_parts={K1 S1 I0 C0 F0}, diff_note={concepts: Brewster angle | steps: refractive index from tan of the angle; wavelength in the medium}]
A 750-nm light beam hits the flat surface of a certain liquid, and the beam is split into a reflected ray and a refracted ray. If the reflected ray is completely polarized and angle of refraction is $37^\circ$, what is the wavelength of the refracted ray ?
\begin{qoptions}
\opt{A}{1000 nm}
\opt{B}{562.5 nm}
\opt{C}{500 nm}
\opt{D}{1125 nm}
\end{qoptions}
\end{bankq}
\begin{bankq}[id=TOAS17-P2-PHY-S1-Q02, source={TOAS - Test 17 (18 Sep 2022)}, paper=2, qno=2, type=single, subtopic={Interference}, answer={C}, answer_src=key, difficulty=1, diff_parts={K1 S1 I0 C0 F0}, diff_note={concepts: intensity in an interference pattern | steps: phase difference from the path difference; cos squared}]
The path difference between two interfering waves at a point on the screen is $\lambda/6$. The ratio of intensity at this point and that at the central bright fringe will be : (Assume that intensity due to each slit in same)
\begin{qoptions}
\opt{A}{0.853}
\opt{B}{8.53}
\opt{C}{0.75}
\opt{D}{7.5}
\end{qoptions}
\end{bankq}
\begin{bankq}[id=TOAS17-P2-PHY-S1-Q09, source={TOAS - Test 17 (18 Sep 2022)}, paper=2, qno=9, type=single, subtopic={Diffraction}, answer={A}, answer_src=key, difficulty=4, diff_parts={K2 S2 I1 C0 F0}, diff_note={concepts: single slit, double slit patterns | steps: identify the pattern type; fringe spacing or width from the picture; compute the width or separation}]
Laser light with a wavelength of 500 nm is directed through one slit or two slits and allowed to fall on a screen 2.60 m beyond. Figure shows the pattern on the screen, with a centimeter ruler below it. Did the light pass through one slit or two slits ? If one slit, find its width. If two slits, find the distance between their centers.
\bfig{148.7}{TOAS17_p2_phy_s1_q09_1.png}
\begin{qoptions}
\opt{A}{One slit of width approx 0.1 mm.}
\opt{B}{One slit of width approx 0.05 mm.}
\opt{C}{Two slits at an approx. distance of 0.08 mm}
\opt{D}{Two slits at an approx. distance of 0.04 mm}
\end{qoptions}
\end{bankq}
\begin{bankq}[id=TOAS17-P2-PHY-S1-Q14, source={TOAS - Test 17 (18 Sep 2022)}, paper=2, qno=14, type=single, subtopic={Interference}, answer={B}, answer_src=key, difficulty=4, diff_parts={K1 S2 I1 C1 F0}, diff_note={concepts: two coherent sources with a phase lead | steps: path and phase differences at A and B; resultant amplitudes; intensity ratio}]
$S_1$ and $S_2$ are two coherent sources of radiations separated by distance $100.25\ \lambda$, where $\lambda$ is the wave length of radiation. $S_1$ leads $S_2$ in phase by $\dfrac\pi2$. $A$ and $B$ are two points on the line joining $S_1$ and $S_2$ as shown in figure. The ratio of amplitudes of component waves from source $S_1$ and $S_2$ at $A$ and $B$ are in ratio $1:2$. The ratio of intensity at $A$ to that of $B\left(\dfrac{I_A}{I_B}\right)$ is :
\bfig{154.5}{TOAS17_p2_phy_s1_q14_1.png}
\begin{qoptions}
\opt{A}{$\infty$}
\opt{B}{$\dfrac19$}
\opt{C}{0}
\opt{D}{8.1}
\end{qoptions}
\end{bankq}
\begin{bankq}[id=TOAS17-P2-PHY-S2-Q06, source={TOAS - Test 17 (18 Sep 2022)}, paper=2, qno=6, type=numerical, subtopic={Interference}, answer={8}, answer_src=key, difficulty=2, diff_parts={K1 S1 I0 C0 F1}, diff_note={concepts: visibility of fringes | steps: amplitude ratio; visibility formula}]
If the ratio of the intensity of two coherent sources is 4 then the visibility $\dfrac{I_{\max}-I_{\min}}{I_{\max}+I_{\min}}$ of the fringes is $n$. Find the value of $10n$.
\end{bankq}
\begin{bankq}[id=TOAS17-P2-PHY-S2-Q07, source={TOAS - Test 17 (18 Sep 2022)}, paper=2, qno=7, type=numerical, subtopic={YDSE}, answer={3}, answer_src=key, difficulty=4, diff_parts={K1 S2 I1 C0 F1}, diff_note={concepts: maxima on a screen line | steps: path difference at the ends; maximum order; count}]
Figure shows two coherent sources $S_1$ and $S_2$ vibrating in same phase, $AB$ is an wire lying symmetrically at a far distance from the sources $S_1$ and $S_2$. Let $\dfrac\lambda d=10^{-3}$. $\angle BOA=0.12^\circ$. How many bright spots will be seen on the wire, including points $A$ and $B$.
\bfig{107.8}{TOAS17_p2_phy_s2_q07_1.png}
\end{bankq}
\begin{bankq}[id=TOAS18-P1-PHY-S3-Q03, source={TOAS - Test 18 (30 Oct 2022)}, paper=1, qno=3, type=numerical, subtopic={Interference}, answer={0.1}, answer_src=solved, working={Reflection from water adds a half-wave shift, so destructive interference needs $2h\sin\theta=\lambda$ (first minimum): $h=120/(2\times0.6)=100$ m $=0.1$ km}, difficulty=3, diff_parts={K1 S1 I1 C0 F1}, diff_note={concepts: interference of direct and reflected radio waves | steps: phase change on reflection; path difference 2 h sin theta equals lambda}]
Radio waves of wavelength 120 m from a galaxy reach a radio telescope by two separate path as shown. As galaxy rises in east over water, the minimum intensity of destructive interference occurs when galaxy is $\theta=37^\circ$ above the horizon. Find the height of the radio telescope disk above the water (in km).
\bfig{109.4}{TOAS18_p1_phy_s3_q03_1.png}
\end{bankq}
\begin{bankq}[id=TOAS18-P2-PHY-S2-Q10, source={TOAS - Test 18 (30 Oct 2022)}, paper=2, qno=10, type=multi, subtopic={YDSE}, answer={A,B,C,D}, answer_src=key, difficulty=5, diff_parts={K1 S2 I1 C1 F1}, diff_note={concepts: two-source interference at a point | steps: path difference at O; condition for intensity; intensity fraction; nearest minimum}]
Two coherent sources of same intensity in same phase $S_1$ \& $S_2$ of wavelength $\lambda$ are placed as shown. Distance of screen is much larger than $d$.
\bfig{108.3}{TOAS18_p2_phy_s2_q10_1.png}
\begin{qoptions}
\opt{A}{If $d=\dfrac{3\lambda}2$ then these will be minimum intensity at $O$.}
\opt{B}{If $d=3\lambda$ then intensity at $O$ is maximum.}
\opt{C}{If $d=\dfrac{11\lambda}6$ then intensity at $O$ is 75\% of maximum intensity.}
\opt{D}{If $d=\lambda$ then minima nearest to $O$ is at $\theta=60^\circ$.}
\end{qoptions}
\end{bankq}
\begin{bankq}[id=TOAS18-P2-PHY-S3-Q05, source={TOAS - Test 18 (30 Oct 2022)}, paper=2, qno=5, type=numerical, subtopic={YDSE}, answer={1}, answer_src=key, difficulty=5, diff_parts={K1 S2 I1 C1 F1}, diff_note={concepts: YDSE with a falling screen | steps: fringe positions; velocity of the screen; relative velocity}]
If monochromatic light of $\lambda=5000$\AA\ is incident perpendicular on slits $S_1$ \& $S_2$ ($d=1$mm, $D=1$m). If screen is released from rest under gravity then velocity of $2^{\text{nd}}$ maxima w.r.t $1^{\text{st}}$ maxima on screen at $t=2$ sec (in cm/s). ($O$ is taken as origin.)
\bfig{146.5}{TOAS18_p2_phy_s3_q05_1.png}
\end{bankq}

\section*{TMAS papers}
\begin{bankq}[id=TMAS04-P1-PHY-S1-Q01, source={TMAS - Test 04 (19 Nov 2023)}, paper=1, qno=1, type=single, subtopic={Interference}, answer={D}, answer_src=key, difficulty=6, diff_parts={K2 S2 I2 C1 F0}, diff_note={concepts: split lens interference, imaging | steps: image positions by the lens formula; virtual sources separation; fringe region on the screen; check each}]
In a convex lens of focal length $10$ cm, a small portion of length $2$ mm is removed symmetrically and the two parts are joined together. A source is placed symmetrically $5$ cm from composite system and a screen is placed $120$ cm from composite system.
\bfig{159.9}{TMAS04_p1_phy_s1_q01_1.png}
\begin{qoptions}
\opt{A}{No interference fringes can be seen on screen}
\opt{B}{To see interference fringes, screen should be moved towards composite system.}
\opt{C}{To see interference fringes, source should be moved away from composite system.}
\opt{D}{Interference fringes will be obtained in region of width $24$ cm.}
\end{qoptions}
\end{bankq}
\begin{bankq}[id=TMAS04-P1-PHY-S2-Q12, source={TMAS - Test 04 (19 Nov 2023)}, paper=1, qno=12, type=multi, subtopic={YDSE}, answer={B,C,D}, answer_src=key, difficulty=6, diff_parts={K1 S2 I2 C1 F1}, diff_note={concepts: YDSE with off axis source | steps: path differences; maxima order; film shift; check each}]
In given young's double slit experiment a point sources of wavelength $\lambda=500$ nm is placed slightly off the central axis. Point A and B are in front of upper and lower slit respectively. If intensity from upper \& lower slits is $I_0$ each on the screen, then
\bfig{239.2}{TMAS04_p1_phy_s2_q12_1.png}
\begin{qoptions}
\opt{A}{Intensity at A, O and B is $4I_0$, $2I_0$ and $4I_0$ respectively.}
\opt{B}{$70^{\mathrm{th}}$ maxima is formed at B}
\opt{C}{$30^{\mathrm{th}}$ maxima is formed at A}
\opt{D}{If a glass slab of thickness $0.02$ mm and refractive index $\mu=3/2$ is kept in front of lower slit then central maxima is formed exactly at O.}
\end{qoptions}
\end{bankq}
\begin{bankq}[id=TMAS04-P1-PHY-S3-Q01, source={TMAS - Test 04 (19 Nov 2023)}, paper=1, qno=1, type=numerical, subtopic={Diffraction}, answer={440.00}, answer_src=key, difficulty=5, diff_parts={K1 S2 I1 C1 F1}, diff_note={concepts: single slit diffraction | steps: first minimum condition; secondary maximum condition; solve}]
In a single slit diffraction experiment first minimum of red light ($660$ nm) coincides with first maximum of some other wavelength (in nm) $\lambda$. Calculate $\lambda$.
\end{bankq}
\begin{bankq}[id=TMAS04-P1-PHY-S3-Q02, source={TMAS - Test 04 (19 Nov 2023)}, paper=1, qno=2, type=numerical, subtopic={Interference}, answer={99.64}, answer_src=key, difficulty=5, diff_parts={K1 S2 I1 C1 F1}, diff_note={concepts: thin film anti reflection | steps: phase changes; thickness condition; value}]
A glass lens is coated on one side with a thin film of magnesium fluoride ($\mathrm{MgF_2}$) to reduce reflection from the lens surface. The index of refraction of $\mathrm{MgF_2}$ is $1.38$; that of the glass is $1.50$. What is the least coating thickness that eliminates (via interference) the reflections at the middle of the visible spectrum ($\lambda=550$ nm) ? Assume that the light is approximately perpendicular to the lens surface. Express your answer in nm.
\bfig{169.7}{TMAS04_p1_phy_s3_q02_1.png}
\end{bankq}
\begin{bankq}[id=TMAS04-P2-PHY-S1-Q03, source={TMAS - Test 04 (19 Nov 2023)}, paper=2, qno=3, type=multi, subtopic={Interference}, answer={A,C}, answer_src=key, difficulty=6, diff_parts={K1 S2 I2 C1 F1}, diff_note={concepts: thin film interference | steps: phase changes; path difference; conditions}]
A glass plate of refractive index $\mu_3$ is coated with a thin layer of thickness $t$ and refractive index $\mu_2$ ($\mu_2>\mu_3$). Light of wavelength $\lambda$ travelling in air is incident normally on layer. It is partially reflected at the upper and the lower surfaces of layer. Choose correct option(s)
\begin{qoptions}
\opt{A}{The two reflected waves interfere constructively if $t=\dfrac{\left(n+\frac12\right)\lambda}{2\mu_2}$.}
\opt{B}{The two reflected waves interfere constructively if $t=\dfrac{\left(n+\frac12\right)\lambda}{2\mu_3}$.}
\opt{C}{The two reflected waves interfere destructively if $t=\dfrac{n\lambda}{2\mu_2}$}
\opt{D}{The two reflected waves interfere destructively if $t=\dfrac{n\lambda}{2\mu_3}$}
\end{qoptions}
\end{bankq}
\begin{bankq}[id=TMAS04-P2-PHY-S2-Q03, source={TMAS - Test 04 (19 Nov 2023)}, paper=2, qno=3, type=numerical, subtopic={YDSE}, answer={0.75}, answer_src=solved, working={Intensity at the centre varies as $I_0\sin^2(\pi t/30)$ from the graph; at $t=10$ s it is $I_0\sin^2 60^\circ=0.75I_0$}, difficulty=6, diff_parts={K1 S2 I2 C1 F1}, diff_note={concepts: film in front of a slit | steps: phase from the thickness; intensity vs time graph; value}]
In the front of the upper slit of YDSE apparatus, a thin film of a liquid of refractive index $1.40$ is placed. It is a hot day the liquid starts evaporating from the surface. A beam of light of wavelength $560$ nm is incident onto the YDSE apparatus and the intensity $I$ at the centre of the screen is monitored. Figure gives intensity $I$ as a function of time $t$. The intensity changes because of evaporation from the two sides of the film. Assume that the film is flat and has parallel sides. Also assume that the film's thickness decreases at a constant rate. If the maximum intensity is $I_0$, then find $x$ if intensity at $t=10$ sec is $xI_0$.
\bfig{160.7}{TMAS04_p2_phy_s2_q03_1.png}
\end{bankq}
\begin{bankq}[id=TMAS04-P2-PHY-S2-Q05, source={TMAS - Test 04 (19 Nov 2023)}, paper=2, qno=5, type=numerical, subtopic={Interference}, answer={1.17}, answer_src=solved, working={Take B at the origin on the $x$-axis: $\sqrt{x^2+16\lambda^2}-x=3\lambda$ gives $x=7\lambda/6$, $n=7/6\approx1.17$}, difficulty=5, diff_parts={K1 S1 I2 C1 F1}, diff_note={concepts: two source interference maxima | steps: path difference condition; solve distance}]
An interference is observed due to two coherent sources A and B having phase difference zero, separated by a distance $4\lambda$ along the $y$-axis where $\lambda$ is the wavelength of the source. If the distance of $3^{\mathrm{rd}}$ maxima on $x$-axis ($\Delta x=3\lambda$) from B is $n\lambda$. find the value of $n$.
\bfig{58.3}{TMAS04_p2_phy_s2_q05_1.png}
\end{bankq}

\section*{TNAS papers}
\begin{bankq}[id=TNAS05-P2-PHY-S1-Q02, source={TNAS - Test 05 (09 Jul 2023)}, paper=2, qno=2, type=multi, subtopic={Interference}, answer={A,B,C}, answer_src=key, difficulty=5, diff_parts={K1 S2 I1 C1 F1}, diff_note={concepts: interference of sound in an interferometer | steps: path difference between extremes; wavelength; frequency; amplitude ratio from intensities}]
Figure shows an acoustic interferometer used to study interference of sound $S$ is a source of sound and $D$ is the detector of sound. Path $SBD$ can be varied in length but path $SAD$ is fixed. It is found that sound intensity has minimum value of 10 $\mu$W/m$^2$ at one position of $B$ and maximum value of 90 $\mu$W/m$^2$ at a position 1.65 cm from the first. Speed of sound in air is 330 m/s.
\bfig{158.2}{TNAS05_p2_phy_s1_q02_1.png}
\begin{qoptions}
\opt{A}{The frequency of sound can be 5000 Hz.}
\opt{B}{The frequency of sound can be 15000 Hz.}
\opt{C}{The ratio of amplitude of waves from $SAD$ and $SBD$ is 2 : 1.}
\opt{D}{The ratio of amplitude of waves from $SAD$ and $SBD$ is 3 : 1.}
\end{qoptions}
\end{bankq}
\begin{bankq}[id=TNAS05-P2-PHY-S1-Q05, source={TNAS - Test 05 (09 Jul 2023)}, paper=2, qno=5, type=multi, subtopic={Interference}, answer={A,C,D}, answer_src=key, difficulty=8, diff_parts={K1 S3 I2 C2 F1}, diff_note={concepts: two loudspeaker interference with sweeping frequency | steps: path difference to the observer; first minimum condition; speed of sound; later maxima; later minima; compare}]
Two loud speakers, A and B, are mounted at ground level, a distance $2d=50.0$ cm apart, and are connected in parallel to the output of an amplifier, which is fed from a variable frequency oscillator. An observer, also at ground level, stands at a point $R$ that is 1.20 km from the line joining the loud speakers and 1.20 km from the perpendicular bisector of $AB$ at $O$. $OR=D$, Figure.  The oscillator is adjusted so that its frequency, $f$, rises linearly with time, at 10 Hz per sec from $f=0$ at time $t=0$. The sound heard by the observer drops to a minimum for the first time when $t=40\sqrt2$ sec.
\bfig{132.8}{TNAS05_p2_phy_s1_q05_1.png}
\begin{qoptions}
\opt{A}{The speed of sound is approximately 367 m/s.}
\opt{B}{The speed of sound is approximately 320 m/s.}
\opt{C}{The sound heard by the observer becomes maximum at approximately $t=108.5$ sec.}
\opt{D}{The sound heard by the observer becomes minimum for second time at approximately $t=150$ sec.}
\end{qoptions}
\end{bankq}
\begin{bankpassage}[id=TNAS05-P2-PHY-S2-G11-13, source={TNAS - Test 05 (09 Jul 2023)}, questions={11-13}]
Answer Q.11, Q.12 and Q.13 by appropriately matching the information given in the three columns of the following table.\par
In a general YDSE setup, a light of wavelength 600 nm is used. `$d$' represents the distance between two slits, `$D$' represents the distance between screen and slits and `$y$' represents the distance of fringe from central maxima.
\bfig{101.1}{TNAS05_p2_phy_s2_para11_13_1.png}
{\footnotesize\setlength{\tabcolsep}{2pt}\begin{tabular}{|c|p{0.1\linewidth}|c|p{0.1\linewidth}|c|p{0.1\linewidth}|}
\hline
\multicolumn{2}{|c|}{\textbf{Column-I $(d)$}} & \multicolumn{2}{c|}{\textbf{Column-II $(D)$}} & \multicolumn{2}{c|}{\textbf{Column-III $(y)$}}\\ \hline
(I) & 1 mm & (i) & 1 m & (P) & 1 mm\\ \hline
(II) & 1.5 mm & (ii) & 1.5 m & (Q) & 1.2 mm\\ \hline
(III) & 2 mm & (iii) & 2 m & (R) & 1.5 mm\\ \hline
(IV) & 2.5 mm & (iv) & 2.5 m & (S) & 2 mm\\ \hline
\end{tabular}}
\end{bankpassage}
\begin{bankq}[id=TNAS05-P2-PHY-S2-Q11, source={TNAS - Test 05 (09 Jul 2023)}, paper=2, qno=11, type=single, subtopic={YDSE}, answer={B}, answer_src=key, difficulty=4, diff_parts={K1 S2 I1 C1 F0}, diff_note={concepts: maxima in the double slit experiment | steps: fringe width; position of maxima; evaluate each option; match}, group=TNAS05-P2-PHY-S2-G11-13]
Which of the following corresponding to the maxima ?
\begin{qoptions}
\opt{A}{(III) (ii) (P)}
\opt{B}{(III) (i) (Q)}
\opt{C}{(IV) (iii) (P)}
\opt{D}{(I) (ii) (S)}
\end{qoptions}
\end{bankq}
\begin{bankq}[id=TNAS05-P2-PHY-S2-Q12, source={TNAS - Test 05 (09 Jul 2023)}, paper=2, qno=12, type=single, subtopic={YDSE}, answer={C}, answer_src=key, difficulty=4, diff_parts={K1 S2 I1 C1 F0}, diff_note={concepts: minima in the double slit experiment | steps: fringe width; position of minima; evaluate each option; match}, group=TNAS05-P2-PHY-S2-G11-13]
Which of the following corresponding to the minima ?
\begin{qoptions}
\opt{A}{(II) (i) (R)}
\opt{B}{(I) (iii) (Q)}
\opt{C}{(II) (iii) (S)}
\opt{D}{(II) (ii) (S)}
\end{qoptions}
\end{bankq}
\begin{bankq}[id=TNAS05-P2-PHY-S2-Q13, source={TNAS - Test 05 (09 Jul 2023)}, paper=2, qno=13, type=single, subtopic={YDSE}, answer={A}, answer_src=key, difficulty=4, diff_parts={K1 S2 I1 C1 F0}, diff_note={concepts: fifth maximum in the double slit experiment | steps: fringe width; position of fifth maximum; evaluate each option; match}, group=TNAS05-P2-PHY-S2-G11-13]
Which of the following corresponds to $5^{\mathrm{th}}$ maxima ?
\begin{qoptions}
\opt{A}{(III) (i) (R)}
\opt{B}{(I) (iii) (R)}
\opt{C}{(II) (i) (P)}
\opt{D}{(II) (iii) (Q)}
\end{qoptions}
\end{bankq}
\begin{bankq}[id=TNAS06-P1-PHY-S1-Q03, source={TNAS - Test 06 (30 Jul 2023)}, paper=1, qno=3, type=single, subtopic={Polarization}, answer={B}, answer_src=solved, working={A polarising sheet transmits the component parallel to its pass axis and absorbs the perpendicular one}, difficulty=2, diff_parts={K2 S1 I0 C0 F0}, diff_note={concepts: polarisation, action of a polarising sheet | steps: statement about transmitted and absorbed components}]
Choose the CORRECT option for electric field components of unpolarized light passing through polarizing sheet.
\begin{qoptions}
\opt{A}{An electric field component parallel to polarizing direction is absorbed and perpendicular to polarizing direction get transmitted.}
\opt{B}{An electric field component parallel to polarizing direction is transmitted and perpendicular to polarizing direction get absorbed.}
\opt{C}{An electric field component parallel to polarizing direction is rotated through some angle and perpendicular component is transmitted.}
\opt{D}{An electric field component parallel to polarizing direction is transmitted and perpendicular component get rotated through some angle.}
\end{qoptions}
\end{bankq}
\begin{bankq}[id=TNAS06-P1-PHY-S1-Q04, source={TNAS - Test 06 (30 Jul 2023)}, paper=1, qno=4, type=single, subtopic={Diffraction}, answer={A}, answer_src=key, difficulty=4, diff_parts={K1 S2 I1 C1 F0}, diff_note={concepts: diffraction envelope on interference fringes | steps: position of the first minimum; fringe spacing; count the maxima in the envelope; locate each}]
Consider a set of two slits of width $b=5\times10^{-2}$ cm and separated by a distance of $d=0.1$ cm. They are illuminated by light of wavelength 600 nm. If a convex lens of focal length 10 cm is placed beyond the slits and resulting interference pattern is observed on a screen in the focal plane of lens :-
\begin{qoptions}
\opt{A}{Within the 1st diffraction minima, we observe 3 maxima due to interference.}
\opt{B}{Within the 1st diffraction minima, we observe 5 maxima due to interference.}
\opt{C}{A maxima within the 1st diffraction minima is located at $y=0.03$ mm.}
\opt{D}{A maxima within the 1st diffraction minima is located at $y=0.09$ mm.}
\end{qoptions}
\end{bankq}
\begin{bankq}[id=TNAS06-P1-PHY-S2-Q06, source={TNAS - Test 06 (30 Jul 2023)}, paper=1, qno=6, type=multi, subtopic={Polarization}, answer={B,D}, answer_src=key, difficulty=5, diff_parts={K1 S2 I1 C1 F1}, diff_note={concepts: polarised and unpolarised components through a rotating polariser | steps: intensity law with angle; read maximum and minimum from the graph; solve for the components; angle}]
A beam of light traveling horizontally consists of an unpolarized component with intensity $I_0$ and a polarized component with intensity $I_p$. The plane of polarization is oriented at an angle $\theta$ with respect to the vertical. The figure shows the total intensity $I_{\text{total}}$ after the light passes through a polarizer as a function of the angle $\alpha$, that the pass axis of the polarizer makes with respect to the vertical. Identify the correct statement (s).
\bfig{236.6}{TNAS06_p1_phy_s2_q06_1.png}
\begin{qoptions}
\opt{A}{$\theta=125^\circ$}
\opt{B}{$I_p=20$ W/m$^2$}
\opt{C}{$I_0=17.5$ W/m$^2$}
\opt{D}{$I_0=10$ W/m$^2$}
\end{qoptions}
\end{bankq}
\begin{bankq}[id=TNAS06-P1-PHY-S2-Q07, source={TNAS - Test 06 (30 Jul 2023)}, paper=1, qno=7, type=multi, subtopic={Polarization}, answer={A,B,D}, answer_src=key, difficulty=5, diff_parts={K1 S2 I1 C1 F1}, diff_note={concepts: polaroids on both slits of a double slit setup | steps: intensity from each slit; interference only for parallel components; cases of angle; evaluate}]
Two identical polaroid films are pasted on both the slits of YDSE apparatus. $O$ is the centre of the screen. Unpolarized light is incident on the apparatus. The intensity at $O$ is seen to be $I_0$, when no films are pasted. Choose the correct statement(s).
\bfig{84.5}{TNAS06_p1_phy_s2_q07_1.png}
\begin{qoptions}
\opt{A}{If pass axes of both films are parallel to each other, intensity at $O$ will be $\dfrac{I_0}{2}$.}
\opt{B}{If pass axes are perpendicular to each other, intensity at $O$ is $\dfrac{I_0}{4}$.}
\opt{C}{If pass axes make an angle of $60^\circ$ with each other, intensity at $O$ is $\dfrac{I_0(\sqrt3+1)}{8}$.}
\opt{D}{If pass axes make an angle of $60^\circ$ with each other, intensity at $O$ is $\dfrac{3I_0}{8}$.}
\end{qoptions}
\end{bankq}
\begin{bankq}[id=TNAS06-P1-PHY-S3-Q09, source={TNAS - Test 06 (30 Jul 2023)}, paper=1, qno=9, type=single, subtopic={YDSE}, answer={A}, answer_src=solved, working={A source shift or a slab shifts the central fringe down with constant width; water reduces the wavelength so the width decreases; a wider slit raises the central intensity with the width constant}, difficulty=5, diff_parts={K1 S3 I2 C0 F0}, diff_note={concepts: changes in a double slit pattern | steps: effect of the source shift; slab effect; medium effect; slit width effect; match}]
A Young's double slit apparatus is shown in figure :-
\bfig{108.3}{TNAS06_p1_phy_s3_q09_1.png}
{\footnotesize\setlength{\tabcolsep}{2pt}\begin{tabular}{|c|p{0.38\linewidth}|c|p{0.3\linewidth}|}
\hline
\multicolumn{2}{|c|}{\textbf{List-I}} & \multicolumn{2}{c|}{\textbf{List-II}}\\ \hline
(P) & Source $S$ move upwards & (1) & Central fringe moves upwards\\ \hline
(Q) & A thin glass slab is placed in front of slit $S_2$ & (2) & Central fringe moves downward\\ \hline
(R) & Whole apparatus is immersed in water & (3) & Fringe width remains constant\\ \hline
(S) & Diameter of slit $S_2$ is doubled & (4) & Fringe width decreases\\ \hline
 & & (5) & Intensity of central maxima increases\\ \hline
\end{tabular}}\par
\begin{qoptions}
\opt{A}{$P\to2,3;Q\to2,3;R\to4;S\to3,5$}
\opt{B}{$P\to2,3;Q\to2;R\to4;S\to5$}
\opt{C}{$P\to2,3;Q\to3,5;R\to3,4;S\to5$}
\opt{D}{$P\to3;Q\to2,3;R\to5;S\to3,5$}
\end{qoptions}
\end{bankq}
\begin{bankq}[id=TNAS06-P1-PHY-S3-Q10, source={TNAS - Test 06 (30 Jul 2023)}, paper=1, qno=10, type=single, subtopic={YDSE}, answer={C}, answer_src=solved, working={At $t=0$ the sources lie along the line to the screen: zero path difference is at $y=\pm\infty$ and the maximum at $y=0$. At $t=2$ they are along $y$: first maximum $y=\dfrac{\lambda D}{d}=250\ \mu$m, the thousandth is at $\sin\theta=0.25$, $y=0.258$ m, not in the list}, difficulty=8, diff_parts={K2 S3 I2 C2 F0}, diff_note={concepts: moving point sources, path difference at a distant screen | steps: source positions at each time; path difference formula; zeroth and last maxima; first maximum position; thousandth maximum; compare}]
Two point monochromatic sources of light are free to move in a circle of radius 1 mm in anticlockwise direction without any external torque with angular velocity of $\pi/4$ rad/sec. In list I various times are given and in list II the coordinate is given (here zeroth maxima / minima is a point where path difference is zero and last maxima/minima is the point when path difference is maximum). Given wavelength of light used $=5000$ \AA\ and neglect Doppler's effect
\bfig{167.3}{TNAS06_p1_phy_s3_q10_1.png}
{\footnotesize\setlength{\tabcolsep}{2pt}\begin{tabular}{|c|p{0.35\linewidth}|c|p{0.3\linewidth}|}
\hline
\multicolumn{2}{|c|}{\textbf{List-I}} & \multicolumn{2}{c|}{\textbf{List-II}}\\ \hline
(P) & $t=0$, zeroth maxima is at & (1) & (1 m, 0)\\ \hline
(Q) & $t=0$, last maxima is at & (2) & (1 m, $-\infty$)\\ \hline
(R) & $t=2$, first maxima is at & (3) & (1 m, 0.25 m)\\ \hline
(S) & $t=2$, thousandth maxima is at & (4) & (1 m, 250 $\mu$m)\\ \hline
 & & (5) & None of these\\ \hline
\end{tabular}}\par
\begin{qoptions}
\opt{A}{$P\to1;Q\to2;R\to3;S\to5$}
\opt{B}{$P\to1;Q\to2;R\to3;S\to4$}
\opt{C}{$P\to2;Q\to1;R\to4;S\to5$}
\opt{D}{$P\to2;Q\to3;R\to1;S\to5$}
\end{qoptions}
\end{bankq}
\begin{bankq}[id=TNAS06-P1-PHY-S4-Q03, source={TNAS - Test 06 (30 Jul 2023)}, paper=1, qno=3, type=numerical, subtopic={Interference}, answer={4}, answer_src=key, difficulty=9, diff_parts={K2 S3 I2 C2 F1}, diff_note={concepts: biprism fringes, intensity ratio | steps: deviation by each prism; overlap region; intensity at the centre; intensity below the point; ratio}]
The Fresnel biprism is a variation in the YDSE. The fresnel biprism consists of two thin prism joint at their bases. Consider a Fresnal biprism experiment (shown in figure) in which prism angle $2\times10^{-3}$ rad, distance of source (of wavelength 600 nm) from prism is 1m and distance between screen and prism is 4m. Find ratio of maximum intensity on screen to intensity just below point $A$ in interference region. [Assuming intensity of light from upper and lower prism is same in interference region and Refractive index of prism $=1.5$]
\bfig{150.3}{TNAS06_p1_phy_s4_q03_1.png}
\end{bankq}
\begin{bankq}[id=TNAS06-P1-PHY-S4-Q04, source={TNAS - Test 06 (30 Jul 2023)}, paper=1, qno=4, type=numerical, subtopic={Interference}, answer={3}, answer_src=solved, working={Reflected rays from the image through the mirror edges reach heights $2d$ and $d/2$ on the screen at $3D$, so $PQ=\tfrac32d=1.5$ mm and $2\alpha=3$}, difficulty=9, diff_parts={K2 S3 I2 C2 F1}, diff_note={concepts: Lloyd mirror geometry, region of interference | steps: image of the source; reflected rays from the mirror edges; heights on the screen; overlap length; double it}]
Refer to the shown diagram. A plane mirror of length $D$, a point source $S$ and a screen are placed as shown. The length $PQ$ on the screen where interference pattern is observed is $\alpha$ mm. The value of $2\alpha$ is
\bfig{102.7}{TNAS06_p1_phy_s4_q04_1.png}
\end{bankq}
\begin{bankq}[id=TNAS06-P1-PHY-S4-Q05, source={TNAS - Test 06 (30 Jul 2023)}, paper=1, qno=5, type=numerical, subtopic={Interference}, answer={2}, answer_src=solved, working={No phase change at either surface (index falls each time), so destructive reflection needs $2nt=(m+\tfrac12)\lambda$, $\lambda=\dfrac{780}{m+1/2}$ nm; only $520$ nm is visible, $520/260=2$}, difficulty=5, diff_parts={K1 S2 I1 C1 F1}, diff_note={concepts: thin film reflection with different media | steps: phase change at each surface; destructive condition; visible wavelength; ratio}]
A thin film of plastic ($n=1.56$) is 0.25 $\mu$m thick. It is sandwiched between two glass slabs with refractive indices of 1.58 and 1.52, respectively. White light (400 -- 700 nm) is first incident normally on the slab for which $n=1.58$. If $\lambda$ is the visible wavelength (in nm) which is missing in the reflected light, then find the value of $\dfrac\lambda{260}$ in nm.
\bfig{54.6}{TNAS06_p1_phy_s4_q05_1.png}
\end{bankq}
\begin{bankq}[id=TNAS07-P1-PHY-S2-Q03, source={TNAS - Test 07 (20 Aug 2023)}, paper=1, qno=3, type=numerical, subtopic={Interference}, answer={206}, answer_src=key, difficulty=9, diff_parts={K2 S3 I2 C2 F1}, diff_note={concepts: Billet split lens, region of interference | steps: image of the source; two coherent images; overlap condition; geometry; distance}]
Converging lens of diameter of aperture 5 cm with a focal length $F=50$ cm was cut in half and spaced at a distance $d=5$ mm. A point source of light $S$ is located at a distance $a=75$ cm from the lens symmetrically. At what distance $\ell$ (in cm) from the lens will interference begin ? The gap between the halves of the lens is closed by an opaque body.
\end{bankq}
\begin{bankq}[id=TNAS07-P1-PHY-S2-Q04, source={TNAS - Test 07 (20 Aug 2023)}, paper=1, qno=4, type=numerical, subtopic={Interference}, answer={1100}, answer_src=key, difficulty=6, diff_parts={K2 S2 I1 C1 F1}, diff_note={concepts: air wedge fringes, thickness of a fibre | steps: condition for the dark band; thickness of the fifth dark band; diameter}]
A very fine hair-like fiber is stuck between two microscope slides. As a result of this, there is a wedge of air between them. When the slides are illuminated normally (from above) with light of wavelength 550 nm, bright and dark interference bands are formed. The fiber is seen to lie at the position of fifth dark band counting from the common edge. What is the diameter (in nm) of the fiber?
\bfig{92.0}{TNAS07_p1_phy_s2_q04_1.png}
\end{bankq}
\begin{bankq}[id=TNAS07-P1-PHY-S2-Q05, source={TNAS - Test 07 (20 Aug 2023)}, paper=1, qno=5, type=numerical, subtopic={Polarization}, answer={75}, answer_src=key, difficulty=5, diff_parts={K1 S2 I1 C1 F1}, diff_note={concepts: Malus law with a rotated analyser | steps: intensity after the polariser; transmitted by the analyser; ratio after rotation}]
An unpolarized light beam is incident on the polarizer of a polarization experiment and the intensity of light beam emerging from the analyzer is measured as 100 Lumens. Now, if the analyzer is rotated around the horizontal axis (direction of light) by $30^\circ$ in clockwise direction, the intensity of emerging light will be \underline{\hspace{1cm}}
\end{bankq}
\begin{bankq}[id=TNAS07-P2-PHY-S1-Q02, source={TNAS - Test 07 (20 Aug 2023)}, paper=2, qno=2, type=multi, subtopic={Interference}, answer={B,D}, answer_src=key, difficulty=6, diff_parts={K1 S2 I2 C1 F1}, diff_note={concepts: thin film interference with different surrounding media | steps: phase change at each interface; path difference; condition; check each case}]
Consider a film of thickness $L$ as shown in four different case blow. Notice the observation of the film with light falling normally. Mark the correct statement(s): (Take $L=\dfrac{3\lambda}{2}$, where $\lambda$ is wavelength in $\mu=1.4$)
\bfig{295.1}{TNAS07_p2_phy_s1_q02_1.png}

\begin{qoptions}
\opt{A}{For (1) and (2), the reflection at film interfaces causes zero phase difference for two reflected rays.}
\opt{B}{For (2) and (3), the reflection at film interfaces causes a phase difference of $\pi$ for two reflected rays.}
\opt{C}{For (1), the film will appear dark, if it is observed through reflected rays from film interfaces.}
\opt{D}{For (3), the film will appear dark, if it is observed through reflected rays from film interfaces}
\end{qoptions}
\end{bankq}
\begin{bankq}[id=TNAS07-P2-PHY-S1-Q03, source={TNAS - Test 07 (20 Aug 2023)}, paper=2, qno=3, type=multi, subtopic={Interference}, answer={B,C,D}, answer_src=key, difficulty=4, diff_parts={K1 S2 I1 C0 F1}, diff_note={concepts: fringe separation in a wedge | steps: fringe width formula; dependence on angle and index; check actions}]
The following diagram is an arrangement to view the interference pattern produced by a wedge of liquid between a microscope slide and a glass block. The interference pattern obtained is made of equally spaced parallel fringes. The fringe separation (or fringe width) can not be decreased by which of the following actions.
\bfig{205.2}{TNAS07_p2_phy_s1_q03_1.png}
\begin{qoptions}
\opt{A}{by increasing the angle of the liquid wedge.}
\opt{B}{by using a liquid of smaller refractive index.}
\opt{C}{by using a thicker glass block.}
\opt{D}{by using a longer liquid wedge of the same angle}
\end{qoptions}
\end{bankq}
\begin{bankq}[id=TNAS07-P2-PHY-S3-Q05, source={TNAS - Test 07 (20 Aug 2023)}, paper=2, qno=5, type=numerical, subtopic={Resolving power}, answer={30.00}, answer_src=key, difficulty=6, diff_parts={K1 S2 I2 C1 F1}, diff_note={concepts: resolving power of a microscope and magnification | steps: minimum resolvable distance; angular size; eye resolution; magnifying power}]
We use a simple microscope to magnify an object. The microscope has a numerical aperture of $\sin\alpha=0.24$. The object is so small that the resolving power of the microscope is totally employed. If the diameter of the eye's pupil is $d=4.0$ mm, find the minimum magnifying power of the microscope ?
\end{bankq}

\section*{Star Review (Other) papers}
\begin{bankpassage}[id=STAR03-P1-PHY-S3-G5-6, source={Star Review (Other) - Test 03 (04 Dec 2023)}, questions={5-6}]
A point source $S$ emitting light of wavelength $\lambda=600$ nm is placed at a height equal to $600$ nm above a flat reflecting surface $AB$ (see figure). The intensity of the reflected light is $64\%$ of the incident intensity. Interference fringes are observed on a screen placed parallel to the reflecting surface at distance $2\lambda$ from it :-
\bfig{82.0}{STAR03_p1_phy_s3_para5_6_1.png}
\end{bankpassage}
\begin{bankq}[id=STAR03-P1-PHY-S3-Q05, source={Star Review (Other) - Test 03 (04 Dec 2023)}, paper=1, qno=5, type=numerical, subtopic={Interference}, answer={0.01}, answer_src=key, difficulty=6, diff_parts={K1 S2 I2 C1 F1}, diff_note={concepts: intensity ratio of interference with a reflecting surface | steps: amplitude ratio from the reflection; phase of the reflected wave; maximum and minimum intensity; ratio}, group=STAR03-P1-PHY-S3-G5-6]
The ratio of minimum to maximum intensity in the interference pattern observed on screen close to point '$P$' is :-
\end{bankq}
\begin{bankq}[id=STAR03-P1-PHY-S3-Q06, source={Star Review (Other) - Test 03 (04 Dec 2023)}, paper=1, qno=6, type=numerical, subtopic={Interference}, answer={2012.46}, answer_src=key, difficulty=7, diff_parts={K1 S2 I2 C2 F1}, diff_note={concepts: bright ring radius from a point source and its image | steps: image position; path difference at the screen; condition for the first bright ring; solve for the radius}, group=STAR03-P1-PHY-S3-G5-6]
The radius (in nm) of the smallest bright ring formed on screen is :-
\end{bankq}
\begin{bankpassage}[id=STAR03-P2-PHY-S2-G9-10, source={Star Review (Other) - Test 03 (04 Dec 2023)}, questions={9-10}]
A monochromatic beam of light falls on Young's double slit experiment apparatus as shown in figure. A thin sheet of glass is inserted in front of lower slit $S_2$.
\bfig{129.4}{STAR03_p2_phy_s2_para9_10_1.png}
\end{bankpassage}
\begin{bankq}[id=STAR03-P2-PHY-S2-Q09, source={Star Review (Other) - Test 03 (04 Dec 2023)}, paper=2, qno=9, type=single, subtopic={YDSE}, answer={D}, answer_src=key, difficulty=4, diff_parts={K1 S2 I2 C0 F0}, diff_note={concepts: position of the central fringe with a slab | steps: extra path from the slab; shift direction; possible positions}, group=STAR03-P2-PHY-S2-G9-10]
The central bright fringe can be obtained
\begin{qoptions}
\opt{A}{At $O$ only}
\opt{B}{At $O$ or below $O$ only}
\opt{C}{At $O$ or above $O$ only}
\opt{D}{Anywhere on the screen}
\end{qoptions}
\end{bankq}
\begin{bankq}[id=STAR03-P2-PHY-S2-Q10, source={Star Review (Other) - Test 03 (04 Dec 2023)}, paper=2, qno=10, type=single, subtopic={YDSE}, answer={A}, answer_src=key, difficulty=5, diff_parts={K1 S2 I2 C1 F0}, diff_note={concepts: condition for central maximum with a slab and angle | steps: path difference from the slit separation; path from the slab; equate}, group=STAR03-P2-PHY-S2-G9-10]
If central bright fringe is obtained on screen at $O$ :-
\begin{qoptions}
\opt{A}{$(\mu-1)\,t=d\sin\theta$}
\opt{B}{$(\mu-1)\,t=d\cos\theta$}
\opt{C}{$(\mu-1)t+d\sin\theta=0$}
\opt{D}{$\dfrac{t}{\mu-1}=\dfrac{d}{\sin\theta}$}
\end{qoptions}
\end{bankq}

\section*{Nurture (General) papers}
\begin{bankq}[id=NUR01-P2-PHY-S3-Q01, source={Nurture (General) - Test 01 (12 Jun 2021)}, paper=2, qno=1, type=single, subtopic={YDSE}, answer={B}, answer_src=solved, working={Slab adds (2-1.5)t optical path on the S1 side; source side gives mu1 d\^{}2/2D=1 micron; so 0.5t=1 micron, t=2 micron}, difficulty=6, diff_parts={K2 S2 I2 C1 F0}, diff_note={concepts: slab in front of one slit of a double slit, ydse | steps: optical path of the slab; source side path difference; balance}]
In order to get central maxima at the centre of the screen (O), the thickness of slab required is :-
\begin{qoptions}
\opt{A}{1$\mu$m}
\opt{B}{2$\mu$m}
\opt{C}{0.5$\mu$m}
\opt{D}{1.5 $\mu$m}
\end{qoptions}
\end{bankq}
\begin{bankq}[id=NUR01-P2-PHY-S3-Q02, source={Nurture (General) - Test 01 (12 Jun 2021)}, paper=2, qno=2, type=single, subtopic={YDSE}, answer={B}, answer_src=solved, working={For a 1 micron slab net shift 0.5 micron towards S2: y=-0.5e-6x1/(1.5x1e-3)=-1/3 mm (the printed 1 mm is read as 1 micron)}, difficulty=6, diff_parts={K2 S2 I2 C1 F0}, diff_note={concepts: position of the central maximum with a slab, ydse | steps: extra path from the slab; source path difference; fringe shift}]
If the thickness of the slab is selected 1mm, then the position of the central maxima will be ($Y$ -coordinate) Taking O to be origin
\begin{qoptions}
\opt{A}{$\dfrac13$mm}
\opt{B}{$-\dfrac13$mm}
\opt{C}{$\dfrac16$mm}
\opt{D}{$-\dfrac16$mm}
\end{qoptions}
\end{bankq}
\end{document}
